(1) 以基态能量为0，化学势 \(\mu\le0\)，激发粒子数 \(N_{\rm ex}=V\int d^3p/(2\pi\hbar)^3[e^{\beta(p^2/2m-\mu)}-1]^{-1}\)。引入热波长 \(\lambda_T=\sqrt{2\pi\hbar^2/(mk_BT)}\)，可写 \(N_{\rm ex}=V\lambda_T^{-3}g_{3/2}(e^{\beta\mu})\)，其中 \(g_s(z)=\sum_{j\ge1}z^j/j^s\)。

(2) 当 \(\mu\to0^-\)，\(g_{3/2}(1)=\zeta(3/2)\) 有限，所以激发态最多容纳 \(N_{\rm ex,max}=V\zeta(3/2)/\lambda_T^3\)。令其等于总粒子数 \(N\)，得 \(\boxed{T_c=\frac{2\pi\hbar^2}{mk_B}[n/\zeta(3/2)]^{2/3}}\)，其中 \(n=N/V\)。若 \(\mu>0\)，低动量占据分母会变成非物理解甚至负数，故不允许。

(3) 在 \(T<T_c\)，热平衡化学势固定在0，凝聚数 \(N_0=N-N_{\rm ex}\)，故 \(\boxed{N_0/N=1-(T/T_c)^{3/2}}\)。总激发内能 \(U=(3/2)Vk_BT\lambda_T^{-3}\zeta(5/2)\)，因此 \(U\propto T^{5/2}\)，定容热容 \(C_V=(5/2)U/T\propto T^{3/2}\)；忽略基态能量。

(4) 热力学凝聚意味着 \(N,V\to\infty\) 且密度固定时基态仍占宏观比例；有限体系的离散谱会把转变抹平。二维均匀连续理想气体的激发粒子数积分在 \(\mu\to0\) 时低动量对数发散，激发态无有限粒子数上限，故均匀二维系统不存在此三维式有限温度理想气体凝聚。
